The most important result of this project is qualitative before it is quantitative: one codebase successfully trained on 2D image benchmarks and a 3D voxel benchmark with no architecture rewrite between the tasks. The same model family, training stack, and CLI workflow were reused across all workloads, with task-specific behavior determined by configuration. That is the core empirical validation of the N-dimensional implementation.

For CIFAR-10, the original preserved sweep selected trial 3 with a validation accuracy of 0.7344 under the 60-epoch sweep budget. After the data-loader randomness bug was fixed, retraining the selected configuration with the 600-epoch budget produced a much stronger restored checkpoint: validation accuracy 0.9182, validation top-5 accuracy 0.9930, and validation loss 0.3237 at epoch 352. A held-out test recomputation of the restored checkpoint reached 0.8995 top-1 accuracy and 0.9936 top-5 accuracy. The training curves in \autoref{fig:cifar10-training} show a stable long retraining run rather than a metric-only artifact.

For CIFAR-100, the fixed 24-trial sweep selected trial 11 with a validation accuracy of 0.5902 and top-5 validation accuracy of 0.8612 under the 50-epoch sweep budget. Full retraining improved the restored validation checkpoint to 0.6948 top-1 accuracy and 0.8856 top-5 accuracy at epoch 300. The held-out test recomputation is closely aligned with validation: 0.6963 top-1 accuracy, 0.8819 top-5 accuracy, and loss 1.3311. This result is not a state-of-the-art CIFAR-100 claim, but it is a useful sanity check that the repaired 2D pipeline learns a non-trivial 100-class task. The curve is shown in \autoref{fig:cifar100-training}.

For ModelNet40, the stabilized sweep is informative precisely because it no longer exposes invalid configuration failures. All 15 trials completed successfully, and the best sweep configuration, trial 11, achieved a validation accuracy of 0.8061 and a top-5 validation accuracy of 0.9499. Full retraining of that selected configuration did \emph{not} improve on the sweep-best score; instead, the restored validation checkpoint reached 0.7421, with top-5 validation accuracy 0.9306 and validation loss 0.9640. A held-out test recomputation reached 0.6965 top-1 accuracy and 0.9149 top-5 accuracy. This does not imply that the architecture fails on 3D data. Rather, it shows that this particular long retraining schedule did not translate the sweep-best configuration into a stronger final checkpoint.

\begin{figure}[t]
  \centering
  \begin{tikzpicture}
  \begin{groupplot}[
    group style={group size=1 by 2, vertical sep=1.0cm},
    width=0.95\linewidth,
    xmin=1,
    xmax=152,
    tick label style={font=\small},
    label style={font=\small},
    grid style={draw=black!10}
  ]
    \nextgroupplot[
      height=0.24\textheight,
      ylabel={Accuracy},
      ymin=0.28,
      ymax=0.96,
      ymajorgrids=true,
      legend style={draw=none, fill=none, font=\footnotesize, at={(0.5,1.16)}, anchor=south, legend columns=3}
    ]
      \addplot[color=ndsblue, line width=1.0pt] table[col sep=comma, x=epoch, y=train_accuracy] {data/cifar10_training.csv};
      \addlegendentry{Train}

      \addplot[color=ndsgold, line width=1.0pt] table[col sep=comma, x=epoch, y=val_accuracy] {data/cifar10_training.csv};
      \addlegendentry{Validation}

      \addplot[only marks, mark=star, mark size=4pt, color=ndsred, fill=ndsred] coordinates {(112,0.7334000468254089)};
      \addlegendentry{Best checkpoint}

      \addplot[color=ndsred, dashed, line width=0.9pt] coordinates {(112,0.28) (112,0.96)};

    \nextgroupplot[
      height=0.22\textheight,
      xlabel={Epoch},
      ylabel={Loss},
      ymin=0.55,
      ymax=2.80,
      ymajorgrids=true
    ]
      \addplot[color=ndsblue, line width=1.0pt] table[col sep=comma, x=epoch, y=train_loss] {data/cifar10_training.csv};
      \addplot[color=ndsgold, line width=1.0pt] table[col sep=comma, x=epoch, y=val_loss] {data/cifar10_training.csv};
      \addplot[color=ndsred, dashed, line width=0.9pt] coordinates {(112,0.55) (112,2.80)};
  \end{groupplot}
\end{tikzpicture}

  \caption{Training history for the final repaired CIFAR-10 retraining run. The dashed marker highlights the best restored checkpoint at epoch 352 with validation accuracy 0.9182; early stopping terminated the run at epoch 392.}
  \label{fig:cifar10-training}
\end{figure}

\begin{figure}[t]
  \centering
  \begin{tikzpicture}
  \begin{groupplot}[
    group style={group size=1 by 2, vertical sep=1.0cm},
    width=0.95\linewidth,
    xmin=1,
    xmax=350,
    tick label style={font=\small},
    label style={font=\small},
    grid style={draw=black!10}
  ]
    \nextgroupplot[
      height=0.24\textheight,
      ylabel={Accuracy},
      ymin=0.02,
      ymax=0.96,
      ymajorgrids=true,
      legend style={draw=none, fill=none, font=\footnotesize, at={(0.5,1.16)}, anchor=south, legend columns=3}
    ]
      \addplot[color=ndsblue, line width=1.0pt] table[col sep=comma, x=epoch, y=train_accuracy] {data/cifar100_training.csv};
      \addlegendentry{Train}

      \addplot[color=ndsgold, line width=1.0pt] table[col sep=comma, x=epoch, y=val_accuracy] {data/cifar100_training.csv};
      \addlegendentry{Validation}

      \addplot[only marks, mark=star, mark size=4pt, color=ndsred, fill=ndsred] coordinates {(300,0.6948000192642212)};
      \addlegendentry{Best checkpoint}

      \addplot[color=ndsred, dashed, line width=0.9pt] coordinates {(300,0.02) (300,0.96)};

    \nextgroupplot[
      height=0.22\textheight,
      xlabel={Epoch},
      ylabel={Loss},
      ymin=0.75,
      ymax=4.60,
      ymajorgrids=true
    ]
      \addplot[color=ndsblue, line width=1.0pt] table[col sep=comma, x=epoch, y=train_loss] {data/cifar100_training.csv};
      \addplot[color=ndsgold, line width=1.0pt] table[col sep=comma, x=epoch, y=val_loss] {data/cifar100_training.csv};
      \addplot[color=ndsred, dashed, line width=0.9pt] coordinates {(300,0.75) (300,4.60)};
  \end{groupplot}
\end{tikzpicture}

  \caption{Training history for the final repaired CIFAR-100 retraining run. The dashed marker highlights the best restored checkpoint at epoch 300 with validation accuracy 0.6948; early stopping terminated the run at epoch 350.}
  \label{fig:cifar100-training}
\end{figure}

\begin{figure}[t]
  \centering
  \begin{tikzpicture}
  \begin{groupplot}[
    group style={group size=1 by 2, vertical sep=1.0cm},
    width=0.95\linewidth,
    xmin=1,
    xmax=26,
    tick label style={font=\small},
    label style={font=\small},
    grid style={draw=black!10}
  ]
    \nextgroupplot[
      height=0.24\textheight,
      ylabel={Accuracy},
      ymin=0.15,
      ymax=1.02,
      ymajorgrids=true,
      legend style={draw=none, fill=none, font=\footnotesize, at={(0.5,1.16)}, anchor=south, legend columns=3}
    ]
      \addplot[color=ndsblue, line width=1.0pt] table[col sep=comma, x=epoch, y=train_accuracy] {data/modelnet40_training.csv};
      \addlegendentry{Train}

      \addplot[color=ndsgold, line width=1.0pt] table[col sep=comma, x=epoch, y=val_accuracy] {data/modelnet40_training.csv};
      \addlegendentry{Validation}

      \addplot[only marks, mark=star, mark size=4pt, color=ndsred, fill=ndsred] coordinates {(6,0.7572879195213318)};
      \addlegendentry{Best checkpoint}

      \addplot[color=ndsred, dashed, line width=0.9pt] coordinates {(6,0.15) (6,1.02)};

    \nextgroupplot[
      height=0.22\textheight,
      xlabel={Epoch},
      ylabel={Loss},
      ymin=0.0,
      ymax=3.8,
      ymajorgrids=true
    ]
      \addplot[color=ndsblue, line width=1.0pt] table[col sep=comma, x=epoch, y=train_loss] {data/modelnet40_training.csv};
      \addplot[color=ndsgold, line width=1.0pt] table[col sep=comma, x=epoch, y=val_loss] {data/modelnet40_training.csv};
      \addplot[color=ndsred, dashed, line width=0.9pt] coordinates {(6,0.0) (6,3.8)};
  \end{groupplot}
\end{tikzpicture}

  \caption{Training history for the final ModelNet40 retraining run. The dashed marker highlights the best restored checkpoint at epoch 7 with validation accuracy 0.7421; early stopping stopped the run at epoch 27.}
  \label{fig:modelnet-training}
\end{figure}

The comparison between sweep-best and final restored-checkpoint performance is summarized in \autoref{tab:summary-results}. The distinction between sweep-best, validation-selected final checkpoint, and recomputed held-out test metrics remains explicit. No error bars are shown because the preserved evidence for this submission consists of single retained runs rather than repeated-seed experiments.

\begin{table}[t]
  \centering
  \scriptsize
  \caption{Summary of the preserved final results. Sweep metrics correspond to the best completed sweep trial. Validation metrics correspond to the restored checkpoint selected during full retraining. Test metrics were recomputed after training from the same restored checkpoint.}
  \label{tab:summary-results}
  \begin{tabular*}{\linewidth}{@{\extracolsep{\fill}}lrrrrrrr@{}}
    \toprule
    \textbf{Dataset} &
    {\textbf{Sweep Val}} &
    {\textbf{Final Val}} &
    {\textbf{Test}} &
    {\textbf{Test Top-5}} &
    {\textbf{Params}} &
    {\textbf{Time (s)}} &
    {\textbf{Best Ep.}} \\
    \midrule
    CIFAR-10 & \(0.7344\) & \(0.9182\) & \(0.8995\) & \(0.9936\) & \(48{,}809{,}188\) & \(14{,}866.06\) & \(352\) \\
    CIFAR-100 & \(0.5902\) & \(0.6948\) & \(0.6963\) & \(0.8819\) & \(48{,}876{,}478\) & \(12{,}043.98\) & \(300\) \\
    ModelNet40 & \(0.8061\) & \(0.7421\) & \(0.6965\) & \(0.9149\) & \(16{,}423{,}126\) & \(990.99\) & \(7\) \\
    \bottomrule
  \end{tabular*}
\end{table}

Overall, the results support a restrained but clear interpretation. NDSwin-JAX is not yet supported by the repeated, standardized benchmark campaign that would justify state-of-the-art claims. Nevertheless, the preserved final artifact set shows that the implementation is functional, configurable, and reusable across dimensionalities. The final checkpoints are published as Hugging Face model repositories for CIFAR-10 (\url{https://huggingface.co/volodymyr-yelisieiev/ndswin-cifar10-final}), CIFAR-100 (\url{https://huggingface.co/volodymyr-yelisieiev/ndswin-cifar100-final}), and ModelNet40 (\url{https://huggingface.co/volodymyr-yelisieiev/ndswin-modelnet40-final}); \autoref{app:artifacts} lists these links together with the committed metric files from which the tables and plots are regenerated. The project succeeds as a practical work in applied machine learning engineering: it translates a known architecture into a broader implementation setting and validates that translation with credible, traceable experiments.
